% Transport solver benchmark -- generated tables/figures under results/transport/

\subsection{Protocol}

A single player, production and combat disabled, must redistribute troops from
an initial to a target distribution on six preset instances (star, path,
bipartite, and three Poisson-map scenarios). Per tick we record the $L^1$ loss
to target; scores are \emph{ticks to converge} (first tick with loss $<1$) and
the \emph{loss integral} $\sum_t \mathrm{loss}(t)$ --- lower means troops
spent less time out of place, i.e.\ the average-arrival-time criterion that
motivates the Kantorovich objective
\citep{kantorovich1942,villani2009,peyre2019}.

Two solvers compete. \emph{Greedy} descends a potential field obtained by
solving a graph Poisson equation on the deficit --- the dynamical,
Benamou--Brenier-flavored view of the problem \citep{benamou2000}. \emph{OT}
solves a min-cost flow \citep{dantzig1963,orlin1997} with
successive-shortest-path/Johnson potentials \citep{johnson1977} and production
saturation via Frank--Wolfe \citep{frankwolfe1956}, the fixed-point scheme
first used for traffic assignment \citep{leblanc1975}. Load-dependent transit
(bigger groups move slower) places the true problem in the flows-over-time
regime of \citet{koehler2002}, which neither solver models exactly --- a gap
the case study below makes visible.

\subsection{Results}

\begin{figure}[ht]
\centering
\includegraphics[width=\textwidth]{\resultsfig{transport/figures/loss_curves.pdf}}
\caption{$L^1$ loss over time per preset (log scale), greedy vs.\ OT. Note the
three regimes: sharp full convergence (star, bipartite/OT), separated tails
(both distributed Poisson scenarios), and the two anomalies --- the identical
staircase on \code{path\_end} and the identical overlap on
\code{capital\_rush}.}
\label{fig:transport-curves}
\end{figure}

\begin{table}[ht]
\centering
\resultstable{transport/tables/transport}
\caption{Ticks to converge (``--'' = never within 800 ticks) and loss
integral. Lower is better.}
\label{tab:transport}
\end{table}

\subsection{Reading}

OT wins where routing structure matters. On the star it converges in nearly
half the time (136 vs.\ 257 ticks; the greedy curve shows a visible plateau
around loss 40 before its final drop). On the bipartite split OT falls off a
cliff to zero at tick 382 while greedy strands a residue of $\sim$34 troops
indefinitely. On the two distributed Poisson scenarios the curves separate
after $\sim$tick 400 and OT's tail accelerates away --- a 7--9\% lower loss
integral and a $2$--$3\times$ lower final loss.

The two anomalies are as informative as the wins. On \code{capital\_rush}
(everything to one sink) the curves overlap exactly: with a single
destination, gradient descent \emph{is} shortest-path routing, so the OT
machinery buys nothing --- a useful sanity check that the solvers agree when
the problem is trivial. On \code{path\_end} both solvers produce the
\emph{identical} staircase --- loss drops in $\sim$200-unit steps roughly
every 150 ticks, then freezes at exactly 800 --- which indicts the
environment, not the solvers. The stepped shape is consistent with the mass
moving as one large, therefore slow, group down the path (speed
$\propto 1/\text{count}$): each step is the blob crossing one edge, and the
freeze is it becoming effectively stationary --- precisely the load-dependent
transit-time effect of \citet{koehler2002}. Confirming this with
\code{gym\_transport\_diag} is \S\ref{sec:open}, item 4.
